\section{Limitations and conclusion}
\label{sec:limitations}

Four limitations should frame any use of the model, and each is disclosed in the codebase as well as here.

\paragraph{The data edge is frozen at 2024Q2.} The estimation sample ends in 2023Q2 and the conditioning data in 2024Q2, matching the source paper's vintage. Forecasts and shock readings therefore describe the economy as of mid-2024; the out-of-sample comparison in Section~\ref{sec:results} is possible precisely because the edge is frozen, but operational use requires re-running the data pipeline on current vintages, and results will shift with revisions. \emph{Update (July 2026):} the production configuration has since been re-estimated on the extended 1992Q1--2025Q1 sample with data through 2026Q1, and the headline FEVD shares quoted in the abstract and Section~\ref{sec:validation} (37.4 and 42.3 per cent) come from that extended run; the frozen-2024Q2 out-of-sample exercise and the original-sample replication figures are retained unchanged.

\paragraph{Internal Bank series are approximated.} The UK-trade-weighted world GDP and world CPI series are unpublished; the replication's OECD/IMF-based proxies, the assumed lag length, the default prior hyperparameters and the simplified pandemic-prior treatment mean the global block matches the paper qualitatively, not exactly. Two substantive discrepancies stand: the global share of one-year UK CPI variance is about eight points below the paper's approximate 50 per cent, and UK monetary policy is not the largest domestic contributor to CPI variance. Both are plausibly attributable to these approximations, but neither has been traced to a cause, and the posterior band on the first is wider than the gap; they are recorded as unresolved rather than explained.

\paragraph{Forecast skill is not established.} The frozen-origin exercise below is one forecast origin. The 49-origin rolling evaluation (Section~\ref{sec:validation}) finds no accuracy gain over a random walk with drift on any of the eight variables at any of the eight horizons once the 64 tests are adjusted together, and finds the predictive bands under-covering. Nothing in this paper supports a claim that the model forecasts better than a naive benchmark.

\paragraph{This is not an official Bank product.} The replication is an independent open-source reconstruction from the published paper. It is not endorsed by, and does not speak for, the Bank of England; where the replication and the paper disagree, the paper is ground truth. Nor does the model score reforms: within PolicyEngine Macro it is the baseline and conditioning member, and reform scoring belongs to the suite's other models.

\paragraph{Extensions.} Three extensions follow naturally from the machinery already in place. Conditional forecasting in the sense of \citet{waggoner1999} --- conditioning the fan charts on a market-implied Bank Rate path or an energy-price assumption --- requires only the linear-restriction algebra of Section~\ref{sec:forecasting-math} plus an economic choice of which shocks implement the conditions. Narrative sign restrictions \citep{antolindiaz2018} on the clearly signed 2008 and 2022 episodes in the estimated shock series would sharpen the identified set. And re-running the data pipeline on current vintages would move the data edge forward, turning the out-of-sample exercise of Section~\ref{sec:results} into a live forecast evaluation loop.

\medskip

\noindent Within those limits, the exercise shows that a central bank's operational structural VAR can be rebuilt in the open: the identification implemented and verified draw by draw, the published magnitude for GDP reproduced and the shortfall on CPI quantified with its posterior band rather than renormalised away, out-of-sample performance on the record including where it is unflattering, and the whole apparatus callable through hosted tools. Replication of policy-institution models, with numbers attached, is both feasible and useful --- and this paper is offered as a template for doing it.
